\makeabstract
    {
    Functional Connectivity Analysis During Transitions Between Autonomous and Manual Driving: An wPLI-based EEG Study
    }
    {
    Ren Sato\textsuperscript{1}, Yue Hu\textsuperscript{2}, Youngbin Kwak\textsuperscript{2}
    }
    {
    \textsuperscript{1} Amherst College, Amherst, MA\\
\textsuperscript{2} University of Massachusetts, Amherst, MA
    }
    {
    Automated driving systems (ADS) at SAE Level 3 allow drivers to disengage from active road monitoring until intervention is needed. While prior studies tracked state shifts using spectral power alone, electrode-level power cannot capture network interactions and is vulnerable to volume-conduction artifacts. This study uses the weighted Phase Lag Index (wPLI) to assess functional connectivity and map neural network shifts during transitions between autonomous (AV) and manual driving (MV). We hypothesized that MV demands greater active cognitive resources and motor control, leading to increased wPLI connectivity in both Theta (4{\textendash}8 Hz) and Alpha (8{\textendash}12 Hz) bands,
    }
    {
    Continuous 32-channel EEG was collected from 28 participants (mean age = 19.90; 60.7\% male) in a simulator. Participants drove 5 scenarios (bicyclist, T-junction, fog, pedestrian, intersection) in AV and MV modes. Data were epoched within scenario boundaries and matched between conditions. Spectral analysis was performed via MTMFFT in FieldTrip to extract Theta (4{\textendash}8 Hz) and Alpha (8{\textendash}12 Hz) bands. Differences in wPLI between conditions were assessed via paired t-tests with False Discovery Rate (FDR) correction.
    }
    {
    Theta Band (4{\textendash}8 Hz): 17 significant edges showed greater connectivity in MV than AV. This increase was heavily anchored at the central sensorimotor hub (Cz), connecting to frontal, temporal, and parietal channels (e.g., F8{\textendash}Cz, Cz{\textendash}TP10, Cz{\textendash}FT9, T8{\textendash}Cz; p \ensuremath{<} 0.001). Alpha Band (8{\textendash}12 Hz): Connectivity remained largely stable, showing only 1 significant edge (FC2{\textendash}CP5; p \ensuremath{<} 0.001) elevated in MV.
    }
    {
    Transitioning from autonomous to manual driving does not universally increase brain connectivity. Instead, manual control specifically recruits a structured, Cz-centered Theta-band sensorimotor network for active steering and cognitive control. Alpha stability suggests drivers maintain visual alertness in both conditions, but physical control requires engaging a distinct higher-order network.
    }

    \newpage
    